\scanpage{6}
120\hfill J.-P. \textsc{Serre}

[1], tome II, p. 177, th. 15) que $\mathfrak s$ est \textit{algébrique} (i.e. correspond à un groupe algébrique); si $\mathfrak c_{\mathrm{alg}}$ (resp. $\mathfrak g_{\mathrm{alg}}$) désigne la plus petite sous-algèbre de Lie algébrique de $\operatorname{End}(V)$ contenant $\mathfrak c$ (resp. contenant $\mathfrak g$), on a évidemment
\[
\mathfrak g_{\mathrm{alg}}=\mathfrak c_{\mathrm{alg}}\times\mathfrak s.
\]
L'algèbre $\mathfrak c_{\mathrm{alg}}$ est abélienne, et opère de façon semi-simple sur $V$. Il s'ensuit que $\mathfrak g_{\mathrm{alg}}$ est réductive. Le lien entre $\mathfrak g_{\mathrm{alg}}$ et $G_{\mathrm{alg}}$ est fourni par la proposition suivante:

\textbf{Proposition 2.} -- (i) \textit{L'algèbre de Lie de $G_{\mathrm{alg}}$ est $\mathfrak g_{\mathrm{alg}}$.}

(ii) \textit{$G_{\mathrm{alg}}$ est réductif} (\textsc{Mumford} [4], p. 26, déf. 1.4).

(iii) \textit{Toute composante connexe de $G_{\mathrm{alg}}$ rencontre $G$.}

Puisque $\mathfrak g_{\mathrm{alg}}$ est algébrique, il existe un sous-groupe algébrique connexe $H^0$ de $\mathbf{GL}_V$ d'algèbre de Lie $\mathfrak g_{\mathrm{alg}}$. Comme $\mathfrak g_{\mathrm{alg}}$ contient $\mathfrak g$, le groupe $H^0(\mathbf Q_p)$ contient un sous-groupe ouvert $G_1$ de $G$; on peut supposer $G_1$ distingué dans $G$. Soit $(g_i)$ un ensemble fini de représentants des classes de $G$ mod. $G_1$; la réunion $H$ des $H^0g_i$ est un sous-groupe algébrique de $\mathbf{GL}_V$ contenant $G$, et il est clair que c'est le plus petit possible. On a donc $H=G_{\mathrm{alg}}$, ce qui démontre (i) et (iii). L'assertion (ii) résulte de (i) (ou bien, si l'on préfère, du fait que $V$ est un $G_{\mathrm{alg}}$-module semi-simple et fidèle).

\textbf{Corollaire.} -- \textit{Si $\mathfrak c$ est algébrique, l'algèbre de Lie de $G_{\mathrm{alg}}$ est égale à $\mathfrak g$, et $G$ est un sous-groupe ouvert de $G_{\mathrm{alg}}(\mathbf Q_p)$.}

La première assertion est évidente; la seconde en résulte puisque $G$ et $G_{\mathrm{alg}}(\mathbf Q_p)$ sont des groupes de Lie $p$-adiques de même algèbre de Lie.

On peut donner une caractérisation des points de $G_{\mathrm{alg}}$ au moyen des \textit{invariants} du groupe $G$. Plus précisément, soient $r,s$ deux entiers $\geq0$, et soit $T_{r,s}(V)$ le produit tensoriel de $r$ copies de $V$ et de $s$ copies du dual $V^*$ de $V$. Le groupe $G$ opère (par transport de structure) sur $T_{r,s}(V)$; soit $T^0_{r,s}(V)$ le sous-espace de $T_{r,s}(V)$ formé des éléments invariants par $G$. Il est clair que $G_{\mathrm{alg}}$ laisse invariants les éléments de $T^0_{r,s}(V)$. De plus, cette propriété \textit{caractérise} les points de $G_{\mathrm{alg}}$:

\textbf{Proposition 3.} -- \textit{Soit $k$ une $\mathbf Q_p$-algèbre commutative, et soit $g\in\mathbf{GL}_V(k)$. Pour que $g$ appartienne à $G_{\mathrm{alg}}(k)$, il faut et il suffit que, pour tout couple $(r,s)$, l'extension $T_{r,s}(g)$ de $g$ à $T_{r,s}(V\otimes k)=T_{r,s}(V)\otimes k$ laisse invariants les éléments de $T^0_{r,s}(V)$.}

Notons d'abord qu'un élément de $T_{r,s}(V)$ est invariant par $G_{\mathrm{alg}}$ si et seulement si il appartient à $T^0_{r,s}(V)$. La prop. 3 est donc conséquence du lemme suivant:

\textbf{Lemme 1.} -- \textit{Soit $E$ un corps de caractéristique 0, soit $V$ un $E$-espace}
\margin{5}
